\section{Full support degree and the original arithmetic quotient}
\label{sec:full-support-degree}
The question addressed here is Deligne's purity condition, with its
actual receiving maps. The supplied English mathematical edition of
Weil II, (3.3.4)--(3.3.6), identifies the image of compactly supported
cohomology in ordinary cohomology as the pure object, after proving
opposite weight bounds and using duality \cite{DeligneII}. We do not
assume an analogous assertion over the split base. We calculate first
what the full support changes in the relative arithmetic group, then
construct an exact map from its prime-degree operator to the original
arithmetic companion action. This includes the attained metric of that
map and its entire distributional extension criterion.

The prime point $\{\tau,e\}$ is retained from the original globalization
paper \cite{globalization}; it is not a discovery of this continuation.
The scalar recollement and prime operators are proved in the preceding
chapters, (RP21)--(RP35), (PD1)--(PD21), and (IA1)--(IA9).
The full finite-support recollement is proved in the accompanying
chapter. The arithmetic matrix is exactly the retained programme
matrix in CLM13--CLM23 of the complete reconstruction source supplied
with this edition. No canonical metric or sheaf-theoretic trace is
inferred from a formal resemblance of operators.

\subsection{The canonical diagonal and all additional support classes}
For a nonempty finite support space $S$, let $c$ be its number of
connected components. The full recollement calculation gives
\[
 D_S=H^1_X(Y,U\mathcal O_X)=\mathbb Q^c/\Delta\mathbb Z,
 \qquad X=\operatorname{Spec}\mathbb Z,
 \tag{FD1}
\]
where $\Delta t=(t,\ldots,t)$ has every component present.
The remaining positive relative degrees are
$H_X^{j+1}(Y,U\mathcal O_X)=H^j(S,\mathbb Q)$ for $j\ge1$.
We shall use the complete group in (FD1), without selecting a branch.
There is a canonical exact sequence
\[
 0\longrightarrow D=\mathbb Q/\mathbb Z
 \xrightarrow{\Delta}D_S
 \xrightarrow{\pi}\mathbb Q^c/\Delta\mathbb Q
 \longrightarrow0.
 \tag{FD2}
\]
The first map is injective because $\Delta q\in\Delta\mathbb Z$
exactly when $q\in\mathbb Z$. Its image is the kernel of the quotient
map $\pi$. That map is onto by the definition of its codomain. This
proves exactness, including the diagonal lattice, rather than merely
an abstract dimension count. Moreover the torsion subgroup of $D_S$
is exactly $\Delta D$: if $n[x]=0$, then $nx\in\Delta\mathbb Z$,
so $x\in\Delta\mathbb Q$; conversely each diagonal rational class
has finite order. This characterization makes the first term canonical.

For a nonzero integer $n$, multiplication $[x]\mapsto[nx]$ is onto,
since $[x/n]$ is a preimage, and
\[
 \ker(n:D_S\to D_S)
   =\Delta(n^{-1}\mathbb Z)/\Delta\mathbb Z,
 \qquad |\ker n|=|n|.
 \tag{FD3}
\]
Every fibre is a translate of this displayed kernel. On
$D_S/\Delta D$ the same multiplication is an automorphism, because
that quotient is a rational vector space. Multiplication by $n$ on
each $H^j(S,\mathbb Q)$, $j\ge1$, is also an automorphism.
Thus disconnected support adds actual classes but does not replace
the arithmetic degree $|n|$ by $|n|^c$. A cycle in the support adds
higher relative cohomology, whose multiplication-by-$n$ degree is one.

For Boolean support with $c=r$ labelled points, forgetting any set of
coordinates, with at least one coordinate retained, gives the actual
map $D_r\to D_t$ by coordinate projection. Its kernel is isomorphic
to $\mathbb Q^{r-t}$: subtract the common integer from the retained
coordinates of a representative whose image vanishes, leaving those
coordinates zero. The deleted coordinates are then uniquely determined
rationals. These maps commute with (FD3), and are the identity on the
diagonal torsion subgroup. This proves their relation to the scalar
branch, including the entire kernel of every deletion.

\subsection{The Hilbert receiver of the entire relative group}
Use counting measure on the countable group $D_S$, and set
\[
 (A_{n,S}f)(d)=f(nd),\quad
 A_{n,S}^*g(u)=\sum_{nd=u}g(d),\quad
 A_{n,S}^*A_{n,S}=|n|I.
 \tag{FD4}
\]
The formulas first hold on finitely supported functions. Collecting the
sum by fibres in (FD3) proves both the adjoint formula and the exact
norm identity $\|A_{n,S}f\|^2=|n|\|f\|^2$; hence the bounded
extensions and their adjoints obey (FD4). They retain all support
components and every diagonal torsion class.

Extension by zero along $\Delta:D\to D_S$ gives an isometric map
\[
 J_S:\ell^2(D)\hookrightarrow\ell^2(D_S),\qquad
 A_{n,S}J_S=J_S A_n,\quad A_{n,S}^*J_S=J_S A_n^*.
 \tag{FD5}
\]
Indeed $nd$ is torsion if and only if $d$ is torsion, since the quotient
in (FD2) is uniquely divisible and torsion free. The fibre above a
torsion element is consequently entirely in $\Delta D$; the fibre
above a nontorsion element is entirely outside it. These observations
prove both identities point by point. The canonical torsion subspace
and its orthogonal complement are therefore both invariant under the
operator and its adjoint. The original vector $\delta_{0}$ belongs to
this torsion subspace. The original scalar Weil receiver (PD14)--(PD18)
is recovered exactly there; the additional classes have not been
assigned invented coefficients in that scalar receiver.

For a fixed prime $p$, put $V_{p,S}=p^{-1/2}A_{p,S}$. It is an
isometry and is pure in the precise operator-theoretic sense
\[
 \bigcap_{m\ge0}V_{p,S}^m\ell^2(D_S)=\{0\}.
 \tag{FD6}
\]
If $f$ belongs to this intersection, then for every $m$ it is constant
on each coset of $\ker p^m$, because it is in the range of pullback by
$p^m$. Thus it is constant on every coset of the infinite subgroup
$\Delta(\mathbb Z[1/p]/\mathbb Z)=\bigcup_m\ker p^m$.
Square summability forces each such constant to vanish. This proves
(FD6). Here ``pure isometry'' is its stated intersection property;
it is not being identified with Deligne's purity theorem.

For completeness an explicit minimal unitary extension follows without
assuming a dilation theorem. Put $W_S=\ker V_{p,S}^*$. The projections
$V^mV^{*m}-V^{m+1}V^{*(m+1)}$ have mutually orthogonal ranges $V^mW_S$.
Their finite sums telescope to $I-V^{M+1}V^{*(M+1)}$. Decreasing
orthogonal projections converge strongly to the projection onto their
intersection: the squared norm differences telescope, so the projected
vectors are Cauchy, and their limit lies in every range. By (FD6) that
limit is zero. Hence
$\ell^2(D_S)=\bigoplus_{m\ge0}V^mW_S$. On
$\widetilde H_S=\bigoplus_{m\in\mathbb Z}W_S$ take the bilateral
shift $\widetilde V$ and embed the nonnegative summands as just
identified. The operator
\[
 \widetilde A_{p,S}=\sqrt p\,\widetilde V,
 \qquad \widetilde A_{p,S}^*\widetilde A_{p,S}
        =\widetilde A_{p,S}\widetilde A_{p,S}^*=pI
 \tag{FD7}
\]
extends $A_{p,S}$. The space $W_S$ is nonzero: otherwise $V$ is a
surjective isometry on a nonzero space, contradicting (FD6).
The bilateral shift has spectrum the unit circle and no eigenvectors.
To verify these claims directly, a unit $w\in W_S$ gives orthonormal
translates $w_m$; the unit vectors
$(2M+1)^{-1/2}\sum_{m=-M}^M\lambda^{-m}w_m$ have residual squared
norm $2/(2M+1)$ for $|\lambda|=1$. Neumann series invert the shift
minus $\lambda$ for $|\lambda|\ne1$. An eigenvector on the circle
would have equal component norms at every integer index and is thus
zero. Scaling proves the full circle $|z|=\sqrt p$ for (FD7).
This extends the proved local weight-one norm identity to every class
of the full relative group (FD1).

The Hilbert construction has a precise relation to additive
complexification, which is a different receiver. On finite-support
functions define
\[
 \kappa_S:\mathbb C^{(D_S)}\longrightarrow D_S\otimes_{\mathbb Z}\mathbb C
       \simeq\mathbb C^c/\Delta\mathbb C,
 \qquad \kappa_S(\delta_d)=d\otimes1.
 \tag{FD8}
\]
The quotient of the free complex vector space by the relations
$\delta_0=0$ and $\delta_{x+y}-\delta_x-\delta_y=0$ has exactly the
universal property of the displayed tensor product: a complex-linear
map out of it is an additive map from $D_S$ to a complex vector space,
and extends uniquely by $d\otimes z\mapsto z f(d)$. This proves both
the map and its complete kernel presentation. The tensor product kills
every torsion element, because its order is invertible in $\mathbb C$.
Using (FD2), or its explicit coordinates after one branch is specified,
then proves the last isomorphism. In particular $\kappa_S\delta_0=0$
and the entire diagonal torsion subspace is killed.

For every positive integer $n$, this exact receiver satisfies
\[
 \kappa_S A_{n,S}=\kappa_S,\qquad
 \kappa_S A_{n,S}^*=n\kappa_S.
 \tag{FD9}
\]
Indeed $A_{n,S}\delta_u=\sum_{nd=u}\delta_d$. Each of its $n$
preimages has tensor value $(u\otimes1)/n$, since the kernel in
(FD3) is torsion. Their sum is $u\otimes1$, proving the first
identity. The adjoint sends $\delta_d$ to $\delta_{nd}$, which proves
the second. Thus the prime pullback on counting functions, the
degree-$p$ operation on the relative group, and its additive
complexification are linked by specified maps with their exact
factors. They are not silently identified. Retaining the supported
zero and torsion is essential for the original integral vector and
the Weil receiver: they disappear under (FD8), whose kernel has now
been calculated rather than treated as absent data.

The supported integer zero is still present. On the full algebraic
dual of $\mathbb C^{(D_S)}\oplus\mathbb C\delta_\tau$, its pullback is
$(f,c)\mapsto(f(0)\mathbf1_{D_S},c)$, whereas absent $\tau$ gives
$(f,c)\mapsto(c\mathbf1_{D_S},c)$. These distinct maps are proved by
evaluating the two scalar actions, precisely as in (IA7). The first
does not preserve $\ell^2(D_S)$, since $D_S$ is infinite. Its exact
target remains the stated algebraic dual; no integer has been deleted.

\subsection{An explicit quotient onto the original arithmetic action}
Keep the original polynomial, its coordinate, and its arithmetic prime:
\[
 \begin{gathered}
 q=(k+1)^2,\quad k\ge17,\quad k\equiv1\pmod4,
 \quad\delta,\gamma>0,\quad p\text{ prime},\\
 Q_k(y)=\prod_{a,b=0}^k[y-(2b-k)\gamma+i(2a-k)\delta],
 \quad E=\mathbb C[y]/(Q_k),\quad M[f]=[yf],\\
 F_p=p^{k/2}\exp(i(\log p)M),\qquad B_p=\exp(i(\log p)M).
 \end{gathered}\tag{AQ1}
\]
The letter $B_p$ denotes this complete centred matrix, not a changed
prime action: $F_p=p^{k/2}B_p$ is retained in every intertwining formula.
At each original cutoff $N\in\{q-1,q,2q-1,2q\}$, let $G_N$ be
the original positive definite metric on $E$ in its original monomial
coordinates. No metric is chosen in its place.

The $k$-fold local relative operator (PD20) acts on
$L^2(\mathbb Q_p^k)$ as $p^k f(x/p)$; its associated unitary operator
is $\mathscr U^{(k)}f(x)=p^{k/2}f(x/p)$. Set
\[
 A=\mathbb Z_p^\times\times\mathbb Z_p^{k-1},\qquad
 h=\mathbf1_A,\qquad h_m=(\mathscr U^{(k)})^mh,
 \quad w_p=\|h_m\|^2=1-p^{-1}.
 \tag{AQ2}
\]
The supports of the $h_m$ are disjoint because the first coordinate
has valuation $m$. Thus their inner products are
$w_p\mathbf1_{m=n}$. The Haar factor $w_p$ is retained, and the
unchanged $k$-fold degree is $p^k$.
Let $\Phi_{p,N}$ be the finite sums $\sum_m h_m\otimes v_m$ in
$L^2(\mathbb Q_p^k)\otimes(E,G_N)$, with the inherited norm. Define
\[
 \Psi_p\Big(\sum_mh_m\otimes v_m\Big)
       =\sum_m B_p^m v_m\in E.
 \tag{AQ3}
\]
Uniqueness of the orthogonal coefficients makes this map well defined;
the sum is finite. It is onto since $h_0\otimes v$ maps to $v$.
For the actual degree operator $T_p=(T_p^{\rm rel})^{\otimes k}\otimes I_E$,
\[
 \Psi_p T_p=F_p\Psi_p,\qquad
 T_p(h_m\otimes v)=p^{k/2}h_{m+1}\otimes v.
 \tag{AQ4}
\]
Both sides of the first identity, on a basis summand, equal
$p^{k/2}B_p^{m+1}v$. This is an exact typed, surjective morphism from
a concrete subspace of the new relative-prime receiver to the original
arithmetic matrix, with its uncentred factor. It does not assert that
the morphism is bounded in the Hilbert norm or is a geometric
compact-to-ordinary cohomology map. Those are mathematical properties
to calculate, rather than inferred identifications.

Its full algebraic kernel is also explicit. Write a finite sequence as
the Laurent polynomial $v(t)=\sum_m t^m v_m$. Then
\[
 \ker\Psi_p=(tI_E-B_p)E[t,t^{-1}].
 \tag{AQ5}
\]
The right side is in the kernel by evaluation. Conversely subtract
$v(B_p)$ from $v(t)$ term by term. For $m>0$,
$t^mI-B_p^m=(tI-B_p)\sum_{j=0}^{m-1}t^{m-1-j}B_p^j$.
For $m=-r<0$ use
$t^{-r}I-B_p^{-r}=-t^{-r}B_p^{-r}(t^rI-B_p^r)$.
These identities hold since $t$ commutes with the matrix and $B_p$
is invertible. Their sum proves (AQ5), including all negative powers.
The support-preserving lift of (AQ3)--(AQ5) on synchronized full-$L$
carriers is $(v,\lambda)\mapsto(\Psi_pv,\lambda)$; its layerwise
kernel is $(\ker\Psi_p)^\uparrow_L$. Direct evaluation of addition,
scalar multiplication and (AQ4) proves these assertions, with the
zero map $(v,\lambda)\mapsto(0,\lambda)$ retained. This is the
retraction category and its supported zero of CR15--CR23; it does not
identify that map with the constant-$\tau$ map on a free semimodule.

\subsection{The attained quotient metric, including every cross term}
Restrict (AQ3) to $-J\le m\le J$, for an integer $J\ge0$. Its
source metric is $w_p\operatorname{diag}(G_N,\ldots,G_N)$.
The minimum squared norm among all preimages of $v\in E$ is
\[
 \min_{\sum_{m=-J}^J B_p^mv_m=v}
      w_p\sum_{m=-J}^Jv_m^*G_Nv_m
  =v^*H_{p,J;N}v,
 \quad
 H_{p,J;N}=w_p\left(
   \sum_{m=-J}^J B_p^mG_N^{-1}(B_p^*)^m\right)^{-1}.
 \tag{AQ6}
\]
To prove the minimum, put $C=\sum_mB_p^mG_N^{-1}(B_p^*)^m$.
It is positive definite because the $m=0$ term is. The candidate
$v_m^0=G_N^{-1}(B_p^*)^m C^{-1}v$ has the required sum. For any
other preimage write $v_m=v_m^0+u_m$, so $\sum_mB_p^mu_m=0$.
The cross term in its squared norm is
$2w_p\operatorname{Re}(v^*C^{-1}\sum_mB_p^mu_m)=0$.
The remaining extra term is nonnegative and vanishes only for $u_m=0$.
The candidate's squared norm is $w_pv^*C^{-1}v$, proving (AQ6).
This is a new, explicitly identified attained quotient metric; it has
not been substituted for any older native conductor metric.

All original roots are distinct for $\delta,\gamma>0$. Let
$\mathsf V_{(a,b),j}=y_{ab}^j$ be their original evaluation matrix,
$0\le j<q$, and set
\[
 y_{ab}=(2b-k)\gamma-i(2a-k)\delta,quad
 \alpha_{ab}=p^{(2a-k)\delta}e^{i(2b-k)\gamma\log p},\quad
 D_N=\mathsf V G_N^{-1}\mathsf V^*.
 \tag{AQ7}
\]
Polynomial interpolation proves $\mathsf V$ invertible and
$\mathsf V B_p\mathsf V^{-1}=\operatorname{diag}(\alpha_{ab})$.
Consequently the complete covariance in (AQ6) satisfies
\[
 (\mathsf V C\mathsf V^*)_{ij}
   =(D_N)_{ij}\,\mathcal D_J(\alpha_i\overline{\alpha_j}),
 \quad
 \mathcal D_J(z)=\sum_{m=-J}^J z^m
 =\begin{cases}
 z^{-J}(1-z^{2J+1})/(1-z),&z\ne1,\\
 2J+1,&z=1.
 \end{cases}
 \tag{AQ8}
\]
Matrix multiplication gives each summand, and the finite geometric sum
gives the last expression. It retains every pair of roots, every
phase resonance, and every original metric entry. In particular the
variance of the $i$th original root functional is exactly
$w_p^{-1}(D_N)_{ii}\mathcal D_J(|\alpha_i|^2)$.

For a rigorous size comparison, define the actual constants
$d_{-,N}=\lambda_{\min}(D_N)>0$ and
$d_{+,N}=\lambda_{\max}(D_N)$. Multiplying
$d_-I\preceq D_N\preceq d_+I$ by each diagonal power, summing, and
inverting positive matrices gives
\[
 \frac{w_p}{d_{+,N}}\mathsf V^*
   \operatorname{diag}(\mathcal D_J(|\alpha_i|^2)^{-1})\mathsf V
 \preceq H_{p,J;N}\preceq
 \frac{w_p}{d_{-,N}}\mathsf V^*
   \operatorname{diag}(\mathcal D_J(|\alpha_i|^2)^{-1})\mathsf V.
 \tag{AQ9}
\]
The upper bound is at most
$w_pd_{-,N}^{-1}p^{-2J\delta}\mathsf V^*\mathsf V$, because $k$
is odd and $|2a-k|\ge1$: choose the endpoint $m=J$ or $m=-J$ in
the positive sum to get
$\mathcal D_J(|\alpha_{ab}|^2)\ge p^{2J|2a-k|\delta}$.
This is a proved exponential collapse of this quotient metric for the
displayed off-line family, with all comparison constants specified.
It is not a conclusion about the existence of actual zeta zeros.

For the four original cutoff signs
$\varepsilon_{q-1}=\varepsilon_q=1$ and
$\varepsilon_{2q-1}=\varepsilon_{2q}=-1$, the exact return of this
new metric is
\[
 \sum_N\varepsilon_N\log\det H_{p,J;N}
  =-\sum_N\varepsilon_N\log\det
   \big[(D_N)_{ij}\mathcal D_J(\alpha_i\overline{\alpha_j})\big]_{ij}.
 \tag{AQ10}
\]
Indeed the exact determinant is
\[
 \det H=\frac{w_p^q|\det\mathsf V|^2}{\det(\mathsf V C\mathsf V^*)}.
\]
Both displayed common factors cancel since $\sum_N\varepsilon_N=0$.
The $D_N$ factors do not
cancel by this argument and have all been retained. Formula (AQ10)
is an exact receiving identity, not an asserted asymptotic evaluation.

\subsection{The precise distribution space selected by the quotient}
The original action now determines which sequences it can receive.
For any invertible matrix $B$ on a finite-dimensional metric space,
consider the already defined finite-sequence map
$\Psi_B((v_m))=\sum_m B^mv_m$. Define the sequence space
\[
 \mathscr S(\mathbb Z;E)=
 \{(v_m):\sup_m(1+|m|)^r\|v_m\|_{G_N}<\infty
                   \text{ for every integer }r\ge0\}.
 \tag{AQ11}
\]
Its topology is given by exactly these seminorms. We have the complete
criterion
\[
 \Psi_B\text{ has a continuous extension to }\mathscr S(\mathbb Z;E)
 \quad\Longleftrightarrow\quad
 |\lambda|=1\text{ for every eigenvalue }\lambda\text{ of }B.
 \tag{AQ12}
\]
Here is a proof including nonsemisimple matrices. In a Jordan basis,
a block is $\lambda I+K$, $K^d=0$. For every integer $m$,
\[
 (\lambda I+K)^m
 =\sum_{j=0}^{d-1}\binom mj\lambda^{m-j}K^j.
 \tag{AQ13}
\]
For nonnegative $m$ this is the binomial formula. For negative $m$ it
follows by multiplying the finite formal inverse series modulo $K^d$,
or inducting downwards from $m=0$ with that inverse. The coefficients
$\binom mj$ are polynomials in $m$ of degree $j$. If all eigenvalues
have modulus one, the finite Jordan change of basis gives an explicit
constant $C_B$ with $\|B^m\|_{G_N}\le C_B(1+|m|)^{d_{\max}-1}$.
Thus the series defining $\Psi_B$ converges absolutely on (AQ11),
bounded by $C_B\sum_m(1+|m|)^{-2}$ times its seminorm of order
$d_{\max}+1$. It is continuous and extends the finite sum.

Conversely continuity of a linear map from (AQ11) to finite-dimensional
$E$ bounds it by $C\sup_m(1+|m|)^r\|v_m\|$ for some finite $C,r$:
continuity at zero supplies finitely many seminorms, and the largest
dominates all of them. Test a sequence supported at the single index
$m$ with coefficient an eigenvector $v$. It gives
$|\lambda|^m\|v\|\le C(1+|m|)^r\|v\|$ for all integers $m$.
Letting $m$ tend to positive and negative infinity proves $|\lambda|=1$.
This also proves uniqueness, since truncation converges in every
Schwartz seminorm (use a seminorm of one higher order for the tail).

For (AQ1), the actual eigenvalues in (AQ7) have moduli
$p^{(2a-k)\delta}$. Thus (AQ12) fails at every $\delta>0$ in the
specified family. The exact root functional at $i=(a,b)$ has norm on
the finite word space
\[
 \|\operatorname{ev}_{y_i}\Psi_p\|^2
   =\frac{(D_N)_{ii}}{w_p}
           \sum_{m=-J}^Jp^{2m(2a-k)\delta}.
 \tag{AQ14}
\]
This is both the diagonal case of (AQ8) and a direct finite dual norm
calculation. It gives the complete rate of its failure to have a
polynomial sequence bound; no projected sign or total determinant
has been used to assign that rate.

At the literal collision $\delta=0$, keep the same polynomial rather
than replacing it by its square-free part:
$Q_k(y)=\prod_{b=0}^k(y-(2b-k)\gamma)^{k+1}$.
The Chinese remainder map gives $k+1$ blocks
$\mathbb C[u]/u^{k+1}$ with $M=(2b-k)\gamma+u$ on the $b$th block.
On that block the original centred powers are
\[
 B_p^m=e^{im(2b-k)\gamma\log p}
       \sum_{j=0}^{k}\frac{(im\log p)^j}{j!}u^j.
 \tag{AQ15}
\]
This follows from the terminating exponential in the same quotient
algebra. Its growth is at most polynomial of degree $k$, so (AQ12)
holds without deleting any repeated root or nilpotent jet. The exact
formula (AQ6) also remains valid at this collision; only the
distinct-root coordinate formula (AQ8) is replaced by (AQ15).

For every parameter, including the collision, the finite-sequence map
has a receiving space selected by its full matrix:
\[
 \mathscr A_B=\{(v_m):\sum_m\|B^m v_m\|_{G_N}<\infty\},
 \qquad\|(v_m)\|_{\mathscr A_B}=\sum_m\|B^m v_m\|_{G_N}.
 \tag{AQ16}
\]
The coordinatewise map $(v_m)\mapsto(B^m v_m)$ is an isometric
bijection to $\ell^1(\mathbb Z;E)$, so this is a complete normed
space. The sum map extends (AQ3) continuously and surjectively to it.
The sequence shift and its inverse are bounded there, with norms at
most $\|B\|$ and $\|B^{-1}\|$, respectively, and the exact
intertwining relation (AQ4) persists. This constructs the object
defined by the failed Schwartz extension, rather than discarding
the arithmetic quotient. The connecting inclusion of finite
sequences into both (AQ11) and (AQ16), and their common map (AQ3),
are explicit. No off-circle eigenvalue has been removed by declaration.

\subsection{The compact-to-ordinary map on the computed affine base}
For the universal coupling on $Y=S\triangleright\operatorname{Spec}\mathbb Z$,
the full recollement proof gives
$R\Gamma(Y,U\mathcal O_X)=R\Gamma(X,\mathcal O_X)$ and
\[
 R\Gamma_X(Y,U\mathcal O_X)
 \longrightarrow R\Gamma(X,\mathcal O_X)
 \longrightarrow R\Gamma(S,\mathbb Q).
 \tag{AQ17}
\]
The middle complex is $\mathbb Z$ in degree zero and zero in positive
degrees. This can be proved directly here: the exact principal-parts
sequence $0\to\mathcal O_X\to\mathbb Q_X\to
\bigoplus_p i_{p*}(\mathbb Q_p/\mathbb Z_p)\to0$ has flasque
last two terms. Constants on nonempty opens of the irreducible $X$
give surjective restrictions in $\mathbb Q_X$; restriction in the
direct sum of point sheaves is deletion of coordinates. The global
principal-parts map is onto by writing a finite collection of
$p$-power fractions and adding them. Its kernel is $\mathbb Z$.
This computes the asserted cohomology, including all primes.

Consequently the natural relative-to-ordinary map in (AQ17) has
zero image in every positive degree, while its boundary groups are
exactly (FD1) and the higher support groups. In degree zero its
relative source is zero, since $\mathbb Z\to\mathbb Q^c$ is injective.
This calculates the particular image; it does not identify a
Zariski support functor with the compact support functor of an
arithmetic curve containing its Archimedean data. The exact
relation is the triangle (AQ17). The relative boundary is retained
as its cone, with the operators (FD3)--(FD7), and is the source
of the arithmetic quotient (AQ3).

These calculations locate the next purity issue precisely. The full
relative spectrum supplies an actual degree operator and its pure
circle norm identity. It also supplies a distributional quotient onto
the original arithmetic action, whose metric and extension space have
now been evaluated. The Zariski relative-to-ordinary image on the
affine universal coupling is zero; it cannot be substituted for the
unconstructed global image involving the Archimedean term and both
pole maps. Establishing that global map's regularity and dual pairing
requires further mathematics. The statements (AQ12)--(AQ16) specify
the exact objects and growth distinction that such a map must address;
they are complete calculations, not an assumed purity theorem or an RH
conclusion.

\begin{figure}[ht]
\centering
\begin{tikzpicture}[node distance=13mm,
 every node/.style={align=center},
 box/.style={draw,rounded corners,text width=118mm,inner sep=5pt}]
\node[box] (coh) {$D_S=\mathbb Q^c/\Delta\mathbb Z$\\
 $0\to\mathbb Q/\mathbb Z\to D_S\to\mathbb Q^c/\Delta\mathbb Q\to0$};
\node[box,below=of coh] (deg) {Every fibre of $d\mapsto pd$ has $p$ elements\\
 $A_{p,S}^*A_{p,S}=pI$; full support retained};
\node[box,below=of deg] (word) {Actual $k$-fold relative operator on prime words\\
 $T_p(h_m\otimes v)=p^{k/2}h_{m+1}\otimes v$};
\node[box,below=of word] (original) {Original quotient $\Psi_p(\sum h_m\otimes v_m)=\sum B_p^mv_m$\\
 $\Psi_pT_p=F_p\Psi_p$; full finite metric (AQ6)--(AQ10)};
\node[box,below=of original] (growth) {Actual root growth $|\alpha_{ab}|=p^{(2a-k)\delta}$\\
 Schwartz extension exactly at unit moduli; repeated jets retained};
\draw[-{Latex}] (coh)--node[right]{(FD3)--(FD4)}(deg);
\draw[-{Latex}] (deg)--node[right]{(PD20), (AQ2)}(word);
\draw[-{Latex}] (word)--node[right]{(AQ3)--(AQ5)}(original);
\draw[-{Latex}] (original)--node[right]{(AQ12)--(AQ16)}(growth);
\end{tikzpicture}
\caption{The exact link from the added prime and full support to the
original arithmetic action. The prime degree fixes its square-root
factor. The arithmetic quotient keeps every original eigenvalue,
coordinate and metric; its growth is calculated rather than assigned.
The diagram does not identify this quotient with Deligne's geometric
image. The latter comparison is (AQ17) and the scope following it;
Deligne's human-source theorem is Weil II, (3.3.6) \cite{DeligneII}.}
\end{figure}
